\documentclass[10pt,a4paper]{article}
\usepackage[utf8]{inputenc}
\usepackage[T1]{fontenc}
\usepackage{lmodern}
\usepackage[dutch]{babel}
\usepackage{amsmath,amssymb,booktabs,longtable,array,geometry,xcolor,hyperref,microtype}
\geometry{left=23mm,right=23mm,top=24mm,bottom=25mm}
\definecolor{SSTblue}{HTML}{173B62}
\definecolor{SSTgray}{HTML}{5F6772}
\hypersetup{colorlinks=true,urlcolor=SSTblue,linkcolor=SSTblue}
\newcommand{\vswirl}{\mathbf{v}_{\!\boldsymbol{\circlearrowleft}}}
\newcommand{\rhof}{\rho_{\!f}}
\newcommand{\rhoE}{\rho_{\!E}}
\newcommand{\rhocore}{\rho_{\mathrm{core}}}
\newcommand{\rc}{r_c}
\newcommand{\status}[1]{\textsf{#1}}
\newcommand{\code}[1]{\texttt{#1}}
\setlength{\parskip}{4pt}
\setlength{\parindent}{0pt}
\renewcommand{\arraystretch}{1.15}
\begin{document}
\begin{center}
{\huge\bfseries\color{SSTblue} SST Mathematical Evidence Registry}\\[5pt]
{\Large v0.1.0 --- Proof-first source and formula audit}\\[7pt]
{\large 20 oktober 2026}\\[9pt]
\fcolorbox{black!35}{black!4}{\parbox{.88\textwidth}{\centering\bfseries DEVELOPMENT / NOT PHYSICALLY CERTIFIED\\[2pt]
41 individueel geauditeerde claims; 379 ongecontroleerde ruwe formuleblokken; 23 regressietests PASS.}}
\end{center}
\vspace{6mm}
\section*{Managementsamenvatting}
Dit register inventariseert de SST-/VAM-literatuur en voorkomt dat eerdere AI-samenvattingen als zelfstandige wetenschappelijke bewijzen worden behandeld.\footnote{Een oudere label als [DERIVED] blijft slechts een \emph{oude beoordeling} totdat een herleidbaar bewijs beschikbaar is.} Het werk bevat (a) bronintegriteit, (b) ongecontroleerde formule-indexering, (c) 41 gekozen formule-audits, (d) kwantitatieve regressietests en (e) een risicoregister. Fysische identiteiten worden niet aangenomen op basis van een wiskundige analogie.

\begin{tabular}{p{.56\textwidth}r}
\toprule
Onderdeel & Resultaat \\
\midrule
SST-ID's in inventaris (00--90) & 91 \\
VAM-ID's in inventaris (01--15) & 15 \\
SST-ID's daadwerkelijk genoemd in de Markdown-samenvatting & 89 \\
Niet genoemd in de Markdown-samenvatting & SST-65, SST-82 \\
SST-ID's met expliciet assessment-heading & 81 \\
Ongecontroleerde \texttt{\$\$...\$\$}-extracties & 379 \\
Geselecteerde, individueel geauditeerde claims & 41 \\
Geregistreerde fouten/ongeldigheden & 10 \\
Lokale regressietests & 23/23 PASS \\
Native PKLSA-kwalificatie & NIET UITGEVOERD \\
\bottomrule
\end{tabular}

\textbf{Interpretatie van de aantallen:} de 379 extracties zijn \emph{geen} 379 bewezen stellingen. Het register heeft voor een groot deel van de SST-00--90-papers slechts de secondaire samenvatting als bron, en de geauditeerde subset kan daarom niet worden geëxtrapoleerd naar de overige papers.

\section{Betrouwbare classificatie}
Voor elke formule worden niet alleen een score, maar een \emph{bron}, \emph{exacte aannames}, \emph{dimensies}, \emph{bewijsniveau}, \emph{numerieke controle} en \emph{falsifier} geregistreerd.
\begin{itemize}\setlength\itemsep{1pt}
\item \status{ESTABLISHED WITH ASSUMPTIONS}: correcte standaardresultaten onder aangegeven voorwaarden; geen SST-bevestiging.
\item \status{CONDITIONAL / MATHEMATICAL CONSTRUCTION}: toetsbaar model met correcte algebra, maar niet als fysica afgeleid.
\item \status{POSTULATED / HYPOTHESIS UNTESTED}: ontbrekende dynamische sluiting of onafhankelijke voorspelling.
\item \status{INCONSISTENT}: aantoonbare fout in eenheid, algebra, numeriek resultaat, topology of variatie.
\item \status{NOT REVIEWED}: ruwe formule uit oudere AI-samenvatting; geen betrouwbaarheidsclaim.
\end{itemize}
De v2-papermatrix (SLV, TRC, FCP enz.) mag als historische \emph{workflowmetadata} worden behouden, maar vervangt geen bewijs op formule-niveau.

\section{Gecontroleerde basis: waar kan de nieuwe SST-actie beginnen?}
Een mechanisch consistente vertrekbasis is het ideale, homogeen-dichte, incompressibele Euler-systeem
\begin{equation}
\rhof(\partial_t\mathbf v+(\mathbf v\cdot\nabla)\mathbf v)=-\nabla p,\qquad \nabla\cdot\mathbf v=0.
\end{equation}
Een materiaalco\"ordinaten-actie met \(J=\det(\partial \mathbf X/\partial\mathbf a)\) is
\begin{equation}
\mathcal A_E=\int dt\int d^3a\,\left[\frac{\rhof}{2}|\partial_t\mathbf X|^2+p(J-1)\right].
\end{equation}
De aparte directorconstructie uit SST-17 kan onder expliciete extra constitutieve aannames worden geschreven als
\begin{equation}
\ell_d=\int d^3x\left[\frac12\rho_d|\partial_t\mathbf q|^2-\frac12 K|\nabla\times\mathbf q|^2\right],\qquad \nabla\cdot\mathbf q=0.
\end{equation}
Hier volgt de vrije golfvergelijking \(\partial_t^2\mathbf q-(K/\rho_d)\nabla^2\mathbf q=0\) (passende randvoorwaarden en \(\rho_d>0,K>0\)). \emph{De koppeling van dit veld aan Euler/PKLSA en de identificatie met fysisch elektromagnetisme zijn niet aangetoond.}

\section{Nieuwe bronnen: twee Kelvin-gap-manuscripten}
De twee verwante Kelvin-manuscripten beschrijven kleine filamentgolven met
\begin{equation}
\omega(k)\simeq\frac{\Gamma k^2}{4\pi}\left[\ln\!\left(\frac1{|k|\xi}\right)+C_0\right], \qquad k_m=\frac{2\pi m}{L}.
\end{equation}
Dit is een bekende \emph{conditionele} dunne-filamentdispersie, geen bewijs voor een topologische energiekloof van een SST-elektron. Voor een \emph{aangenomen} bosonische kloof \(\Delta_K>0\) is de thermische bijdrage
\begin{equation}
U_{\mathrm{th}}=\frac{\Delta_K}{\exp(\Delta_K/k_BT)-1}\sim \Delta_K \exp[-\Delta_K/(k_BT)] \qquad (k_BT\ll\Delta_K).
\end{equation}
De exponenti\"ele onderdrukking is mathematisch correct \emph{als} een kloof bestaat. Het bestaan van \(\Delta_K\) is niet uit Euler of PKLSA bewezen.

\subsection*{Drie onafhankelijke controles}
\textbf{(1) Frequentie--energie:} de eerste paper noemt \(\omega_1\simeq2.8\times10^{17}\,\mathrm{rad/s}\) en \(E_1=29.5\,\mathrm{eV}\). Maar
\begin{equation}
\hbar\omega_1= (1.054571817\times10^{-34}\,\mathrm{J\,s})(2.8\times10^{17}\,\mathrm{s^{-1}})
=\boxed{184.3\,\mathrm{eV}}.
\end{equation}
De beschreven filamentparameters met \(C_0=1\) leveren ongeveer \(2.4565\times10^{17}\,\mathrm{s^{-1}}\), dus \(161.7\,\mathrm{eV}\), niet \(29.5\,\mathrm{eV}\).

\textbf{(2) Circulatie:} het manuscriptaanname \(\Gamma=h/m_e\) en de SST-canon \(\Gamma=2\pi r_c\|\vswirl\|\) zijn niet gelijk:
\begin{align}
\Gamma_{\rm paper}&=7.273895\times10^{-4}\,\mathrm{m^2/s},\\
\Gamma_{\rm canon}&=9.683619\times10^{-9}\,\mathrm{m^2/s},\\
\Gamma_{\rm paper}/\Gamma_{\rm canon}&=\boxed{7.5115\times10^4}.
\end{align}
Aangezien \(\omega(k)\propto\Gamma\), is ook de Kelvin-schaal sterk afhankelijk van het gekozen circulatiecontract. Een vergelijking met \(\Gamma=h/m_e\) kan niet als numerieke SST-canonvoorspelling worden geregistreerd.

\textbf{(3) Dimensies van de voorgestelde gap:} \(\Gamma^2/(4\pi\xi L)\) heeft \(\mathrm{m^2/s^2}\), niet joule. De zogenoemde lijnspanning \(\rho\Gamma^2/(4\pi\xi^2)\) heeft de dimensie van \emph{druk}, niet van lijnspanning. Een afzonderlijk mechanisch afgeleide gapfunctionaal is vereist.

Bij de bijkomende \emph{aanname} \(\Delta_K/(k_BT_{\rm eff})\ge60\) en \(T_{\rm eff}=10^5\,\mathrm K\) volgt rekenkundig \(\Delta_K\ge517.0\,\mathrm{eV}\). De grens van 60 zelf en de spectroscopische bron daarvan zijn nog niet onafhankelijk geverifieerd.

\section{Nieuwe bron: dual-vacuum- en Unruh-echo-ansatz}
De dubbele-substraatpaper bevat de correcte Euler-wervelstrekvergelijking
\begin{equation}
\frac{D\boldsymbol\omega}{Dt}=(\boldsymbol\omega\cdot\nabla)\mathbf v.
\end{equation}
Dit impliceert niet dat \emph{elke} translatie of versnelling een vortex oprekt: de snelheidsgradi\"ent bepaalt de stretching.

De paper definieert \(Z_S=\rho_f\|\vswirl\|\) en \(\kappa_{se}=4Z_S/Z_{\rm bound}\), hetgeen met \(Z_{\rm bound}=10^7\,\mathrm{Rayl}\) oplevert:
\begin{equation}
Z_S=0.76569\,\mathrm{Rayl},\qquad \kappa_{se}=3.063\times10^{-7}\quad\text{(\emph{dimensieloos}).}
\end{equation}
De daaropvolgende populatievergelijking telt echter \(\kappa_{se}\) op bij vervalsnelheden \(\Gamma_S,\gamma_{\rm diss}\) met dimensie \(\mathrm{s^{-1}}\). Dit is ongeldig. Een gerepareerd model vereist een onafhankelijk gedefinieerde overdrachtssnelheid \(k_{SE}\,\,[\mathrm{s^{-1}}]\):
\begin{align}
\dot n_S&=-(\Gamma_S+\gamma_{\rm diss}+k_{SE})n_S,\\
\dot n_{EM}&=k_{SE}n_S-\Gamma_{EM}n_{EM}.
\end{align}
Als \(n_{EM}(0)=0\) en \(\lambda_S=\Gamma_S+\gamma_{\rm diss}+k_{SE}\), dan volgt
\begin{equation}
n_{EM}(t)=\frac{k_{SE}n_S(0)}{\Gamma_{EM}-\lambda_S}\left(e^{-\lambda_St}-e^{-\Gamma_{EM}t}\right).
\end{equation}
\textbf{Status:} correcte oplossing van een tweekompartimentenmodel; geen bewijs dat de twee populaties fysische vacu\"umsectoren zijn of dat een Unruh-echo gemeten is. Een dimensieel onduidelijke vector--directorvermenigvuldiging in de gewijzigde wet van Faraday moet eerst als tensoroperator worden gedefinieerd.

\section{Gekoppelde afhankelijkheids- en probleemlijst}
\small
\begin{longtable}{p{.13\textwidth}p{.30\textwidth}p{.43\textwidth}}
\toprule
Claim & Type & Blokkade / repara\-tie \\
\midrule
\endhead
M001/M006 & massa-schaal & \(\rho v^2V\) is energie; massa vraagt expliciete \(c^{-2}\)-factor of een dimensionale constante. \\
M005 & Hamiltoniaan & Oud VAM-5 gebruikt conjugate impuls voor \(\dot{\mathbf v}\), terwijl \(\dot{\mathbf v}\) in de getoonde actie ontbreekt. \\
M016 & \(0_1\leftrightarrow3_1\) & Knooptype blijft behouden bij gladde ideale Euler-evolutie; reconnectie is een apart mechanisme. \\
M024 & Kelvin numeriek & \(\hbar\omega_1\): 184.3 eV in plaats van 29.5 eV. \\
M025/M026 & Kelvin dimensies & Gap en spanning worden met verkeerde SI-dimensies gepresenteerd. \\
M027/M041 & Kelvin bewijs/constanten & Gap niet afgeleid, en paper-circulatie verschilt \(\sim75{,}115\times\) van canon. \\
M032 & uitwisselingsmodel & Dimensieloos \(\kappa_{se}\) toegevoegd aan vervalsnelheden. \\
M035 & Faraday-koppeling & Multiplicatie vector maal director is geen expliciete vector/tensorkoppeling. \\
\bottomrule
\end{longtable}
\normalsize
Overige niet-gekwalificeerde beweringen zoals fotonidentiteit, Coulomb-veld uit heliciteit, quantumgaps, massahierarchie en zwaartekracht blijven in de catalogus, met een open status; ze worden niet stilzwijgend overgenomen.

\section{De volgende onafhankelijke kwalificatiegates}
\begin{enumerate}\setlength\itemsep{2pt}
\item \textbf{G0--G1:} echte code-/datasetprovenance en concordantie van originele papers met de AI-samenvattingen.
\item \textbf{G2:} parser, eenheden, randvoorwaarden, definities en canonische symbolen.
\item \textbf{G3:} variatierekening, Hamiltoniaanse Legendre-transformatie, Euler/Poisson-sluiting, no-go's.
\item \textbf{G4--G5:} bron-native PKLSA-spectrum; gescheiden \(N=1200/2400\) en \(N=4800/9600\) met ruimtelijke en temporele convergentie; circulatie, energie en heliciteit onafhankelijk meten.
\item \textbf{G6:} blind toetsbare resonantie-/gap-hypothesen tegenover nullmodellen, schaal- en Kelvin-baselines.
\item \textbf{G7:} alleen bij onafhankelijke experimentele correlatie de fysische betekenis promoveren.
\end{enumerate}
Geen van deze gates legitimeert een bestaande filosofie als bewijs. Numerieke regressies betekenen alleen dat de gerapporteerde softwarecontroles slagen.

\section*{Reproduceerbaarheid en bestanden}
De bronkopie\"en en hun SHA-256 hashes staan in \code{results/source\_manifest.json}. \code{build\_registry.py} leest de bronindex en voert de deterministische export naar JSONL/CSV/JSON uit. \code{run\_all.cmd} (Windows) en \code{run\_all.sh} (Linux) starten build en unittest. Elke curatorclaim heeft een claim-ID, papernummer, bron-ID, locatie, LaTeX-formule, dimensiecheck, status en falsifier. De ZIP bevat de originele (ongewijzigde) beschikbare bronbestanden, behalve de redundante PDF-export van de SST-samenvattingen die door de tekst-Markdown vervangen is.

\section*{Referenties}
\begin{thebibliography}{99}
\bibitem{Saffman1992} P.~G. Saffman (1992), \emph{Vortex Dynamics}, Cambridge University Press. \url{https://doi.org/10.1017/CBO9780511624063}.
\bibitem{Moffatt1969} H.~K. Moffatt (1969), ``The degree of knottedness of tangled vortex lines'', \emph{J. Fluid Mech.} \textbf{35}, 117--129. \url{https://doi.org/10.1017/S0022112069000991}.
\bibitem{Unruh1976} W.~G. Unruh (1976), ``Notes on black-hole evaporation'', \emph{Phys. Rev. D} \textbf{14}, 870--892. \url{https://doi.org/10.1103/PhysRevD.14.870}.
\bibitem{Simoncelli2022} M. Simoncelli, N. Marzari and F. Mauri (2022), ``Unified theory of thermal transport in crystals and glasses'', \emph{Nature Physics} \textbf{18}, 1180--1190. \url{https://doi.org/10.1038/s41567-022-01680-2}.
\bibitem{SST29} O. Iskandarani (2025--2026), \emph{Kelvin Mode Suppression in Atomic Orbitals: A Vortex-Filament Gap}, supplied 10-page manuscript, pp.~2--8 (no verified DOI supplied).
\bibitem{SST29B} O. Iskandarani (2025), \emph{Suppression of Kelvin-Mode Thermodynamics in Atomic Orbitals: A Gapped Vortex-Filament Resolution in Swirl--String Theory}, supplied 8-page manuscript, pp.~3--6 (DOI placeholder in source).
\bibitem{DualVac} O. Iskandarani (2025), \emph{Hydrodynamic Dual-Vacuum Unification: Swirl Superradiance, Photon Torsion, and the Unruh Echo Mechanism}, supplied 5-page manuscript, pp.~2--4 (no verified DOI supplied).
\bibitem{SST17} O. Iskandarani (2026), \emph{Electromagnetism as Propagating Torsion in a Hydrodynamic Vacuum}, SST-17 preprint, \url{https://doi.org/10.5281/zenodo.17677074}.
\end{thebibliography}
\end{document}
